\section{Research subjects and scope}

This section defines the core subjects under investigation and delineates the operational boundaries of this study.

\subsection{Research Subjects}

The primary research subjects of this study include:
\begin{itemize}
    \item Graph neural network collaborative filtering architectures for personalized recommendation.
    \item Multimodal feature extraction techniques for clothing items, specifically visual encoders (MobileNetV2 and CLIP) and textual representation models (TF-IDF).
    \item Modal integration and fusion strategies, including offline data-level similarity propagation, online model-level late fusion, and dynamic weighted attention gating.
\end{itemize}

\subsection{Research Scope}

The boundaries of the research are defined as follows:
\begin{itemize}
    \item \textbf{Dataset Scope}: All experiments are conducted on the public Vibrent Clothes Rental (VCR) dataset. Preprocessing is restricted to primary display images and a standard 5-core filter, resulting in a dataset of 553 users, 2,194 items, and 9,455 transactions.
    \item \textbf{Algorithmic Scope}: We evaluate three GNN collaborative filtering architectures (CombiGCN adapted, BM3, and FREEDOM), all built on the LightGCN propagation scheme as their shared collaborative-filtering backbone. We do not investigate sequence-aware recommendation or real-time conversational agents, which are reserved for future work.
    \item \textbf{Feature Scope}: Visual features are extracted using pre-trained, frozen MobileNetV2 and CLIP models, and textual features are extracted using TF-IDF vectors, without fine-tuning the raw encoders.
    \item \textbf{Evaluation Scope}: Performance is measured using six offline top-$N$ ranking metrics (Recall, Precision, NDCG, MRR, MAP, Hit Ratio) evaluated at $K \in \{1, 5, 10, 20\}$.
\end{itemize}
